%\subsection{Schematic Sheets}
% \begin{indentedfigure}
%     \includegraphics[width=1\linewidth]{Images/v0.1/Schematic/Schematic-toplevel.png}
%     \caption{Top-Level schematic sheet from C2Sv0.1 Breakout Board Project. Only critical paths were included on the Top-Level schematic. Routes like general power and ground were not included}
%     \label{fig:v0.1-schematic-toplevel}
% \end{indentedfigure}

% \begin{indentedfigure}
%     \includegraphics[width=1\linewidth]{Images/v0.1/Schematic/C2Sv0-1 Block.png}
%     \caption{General block diagram of C2Sv0.1 breakout board. Components like resistors, capacitors and suggested circuits have been omitted from the diagram. This is meant to show the core functionality between the core components.}
%     \label{fig:c2sv0.1-block-diagram}
% \end{indentedfigure}

% \begin{indentedfigure}
%     \includegraphics[width=1\linewidth]{Images/v0.1/Schematic/Sensor.png}
%     \caption{Image sensor schematic sheet from C2Sv0.1.}
%     \label{fig:v0.1-schematic-sensor}
% \end{indentedfigure}

% \begin{indentedfigure}
%     \includegraphics[width=1\linewidth]{Images/v0.1/Schematic/MCU.png}
%     \caption{Image sensor schematic sheet from C2Sv0.1.}
%     \label{fig:v0.1-schematic-mcu}
% \end{indentedfigure}

% \begin{indentedfigure}
%     \includegraphics[width=1\linewidth]{Images/v0.1/Schematic/Memory.png}
%     \caption{Image sensor schematic sheet from C2Sv0.1.}
%     \label{fig:v0.1-schematic-memory}
% \end{indentedfigure}

% \begin{indentedfigure}
%     \includegraphics[width=1\linewidth]{Images/v0.1/Schematic/USB-C.png}
%     \caption{Image sensor schematic sheet from C2Sv0.1.}
%     \label{fig:v0.1-schematic-usbc}
% \end{indentedfigure}

% \begin{indentedfigure}
%     \includegraphics[width=1\linewidth]{Images/v0.1/Schematic/uSD.png}
%     \caption{Image sensor schematic sheet from C2Sv0.1.}
%     \label{fig:v0.1-schematic-usd}
% \end{indentedfigure}

% \begin{indentedfigure}
%     \includegraphics[width=1\linewidth]{Images/v0.1/Schematic/LPS.png}
%     \caption{Image sensor schematic sheet from C2Sv0.1.}
%     \label{fig:v0.1-schematic-lps}
% \end{indentedfigure}

% \begin{indentedfigure}
%     \includegraphics[width=1\linewidth]{Images/v0.1/Schematic/TAG.png}
%     \caption{Image sensor schematic sheet from C2Sv0.1.}
%     \label{fig:v0.1-schematic-tag}
% \end{indentedfigure}

\subsection{STM32U5 Power Management}

\subsubsection{Power Supplies}\label{sec:app-u5-powersupply}
    The STM32U5 requires various independent supplies for specific peripherals. The following breakdown is a justification of the list from page 3 of the getting started guide.
    \begin{itemize}
        \item \textbf{$V_{DD}$ = 3.3V} \\
        $V_{DD}$ serves as the primary digital power domain for the STM32U5, supplying the processor core's internal regulator, general-purpose I/O banks, and most on-chip peripherals. A 3.3V rail was selected as it represents the most common logic level for mixed-signal embedded systems and is natively supported by all external components in the design, including the SPI flash memory, sensor control lines, and voltage-level interfaces. 
        Selecting 3.3V simplifies the power architecture by allowing a single shared rail across the majority of devices, avoiding the need for intermediate level shifters and reducing regulator count. It also provides greater noise margin and signal integrity on high-speed digital traces compared to lower-voltage domains, at a modest increase in power consumption. The choice therefore balances component compatibility, design simplicity, and system robustness.
    
        \item \textbf{$V_{DDA}$ = 1.8V} \\
        $V_{DDA}$ is the external-analog power supply for A/D converters, D/A converters, voltage reference buffer, operational amplifiers, and comparators \cite{u5-getting-started}. It must be independent from $V_{DD}$ when the mentioned peripherials are used \cite{u5-getting-started}.
        Selecting 1.8V helps reduce the number of different voltage sources, and subsequently number of components, on the board since 1.8V is already a power supply for the image sensor. Moreover, the lower voltage helps reduce overall power consumption.

        \item \textbf{$V_{DDSMPS}$ = 3.3V} \\
        $V_{DDSMPS}$ is the external power supply for the SMPS step-down converter \cite{u5-getting-started}. When in use, it must be connected to the same pin as $V_{DD}$ \cite{u5-getting-started}. So, its voltage is the same as $V_{DD}$, 3.3V.

        \item \textbf{$V_{LXSMPS}$ = $V_{DDCORE}$} \\
        This is the output pin from the internal SMPS step-down converter and has a typical voltage of 1.21V (pg. 166 \cite{u5-datasheet}) \cite{u5-getting-started}.

        \item \textbf{$V_{DD11} = V_{DDCORE}$} \\
        This is the digital core supply for the processor and is supplied by $V_{LXSMPS}$ \cite{u5-getting-started}. It requires a total of $4.7\mu F$ external capacitors, two $2.2\mu F$ bulk capacitors and $100nF$ decoupling capacitor for each pin (in this case) \cite{u5-datasheet} \cite{u5-getting-started}

        \item \textbf{$V_{CAP}$ = UNAVAILABLE} \\
        These pins are not offered on the selected U5G processor.

        \item \textbf{$V_{DDUSB}$ = 3.3V} \\
        This is the external-independent power supply for USB transceivers \cite{u5-getting-started}. Since the USB-C connector is only being used to supply power, the transceiver functionality was not applicable, so, the pin was connected to $V_{DD}$ as specified by \cite{u5-getting-started}.

        \item \textbf{$V_{DD11USB}$ = UNAVAILABLE} \\
        This pin is not offered on the selected U5G processor.

        \item \textbf{$V_{DDIO2}$ = $V_{DDCORE}$} \\
        $V_{DDIO12}$ is the external power supply for 14 I/Os (port G[15:2]) and its voltage level is independent from the $V_{DD}$ voltage \cite{u5-getting-started}. It was connected to $V_{DDA}$ since we require the I2C lines on PG13 and PG14 to run at 3V3 logic. If these I2C lines were not connected, $V_{DDIO2}$ would be connected to $V_{DDCORE}$:
        The 1.8V and 2.8V voltage levels are produced by the local LDOs, whereas the $V_{DDCORE}$ voltage is produced by a much more efficient SMPS step-down converter \cite{stm-power}, which means there are less losses between regulating from SMPS than an LDO. This also allows the pin to still be powered if there are any issues with the 2.8V LDO (if the 1.8V LDO had issues, VREF+ would not be in its operating range and the processor would not function).

        \item \textbf{$V_{BAT}$ = 3.3V} \\
        This is the power supply when $V_DD$ is not present for RTC, TAMP, LSE, backup registers and optionally backup SRAM \cite{u5-getting-started}. The pin has been shorted to $V_DD$ since the camera will not be placed in any situation where it still needs to operate without $V_DD$ (if the main power goes out on the CubeSat of HAB, the camera cannot do anything until the main power comes back).

        \item \textbf{$V_{REF-}$ = 1.8V} \\
        $V_{REF-}$ is the ground reference voltage for the processor. When the VREF+ pin is double-bonded to $V_{DDA}$ in a package, the internal VREFBUF is not available, and must be kept disabled, $V_{REF-}$ must always be equal to $V_{SSA}$ \cite{u5-getting-started}.

        \item \textbf{$V_{REF+}$ = 1.8V} \\
        As outlined in \cite{stm-power}, the simplest solution to provide $V_{REF}$ is to use $V_{DDA}$, 1.8V. Since it is connected to $V_{DDA}$, it admits to the 0.4V voltage difference between $V_{REF}$ and $V_{DDA}$ (pg. 163 \cite{u5-datasheet}).

        \item \textbf{$V_{DDDSI}$ = 3.3V} \\
        This is the external power supply for the DSI controller \cite{u5-getting-started}. It is provided externally through the $V_{DDDSI}$ supply pin, and must be connected to the same supply as $V_{DD}$ pin \cite{u5-getting-started}.

        \item \textbf{$V_{DD11DSI} = V_{DDCORE}$} \\
        This is the external power supply for the DSI transceiver and must be connected to $V_{DD11}$ \cite{u5-getting-started}.
    \end{itemize}
\subsubsection{Decoupling Capacitors}
    The decoupling capacitors followed the getting started guide and datasheet power supply schemes (pg. 14 \cite{u5-getting-started} and pg. 161 \cite{u5-datasheet}).
    The 0603 package was used for all decoupling and capacitors and inductors to comply with the PAST Avionics standard.
     
    \begin{indentedfigure}
        \includegraphics[width=1\linewidth]{Images/Appendices/v0.1//Power Scheme.png}
        \caption{STM32U5GxxxxxQ power supply scheme (with SMPS)}
    \end{indentedfigure}

\subsection{Pin Justification Tables}

% ------------------------
% STM32U5G 144-Pin Justification Tables 
% ------------------------

\begin{table}[H]
    \centering
    \begin{tabular}{|l|p{3cm}|p{6cm}|p{4.5cm}|}
        \hline
        \textbf{Pin} & \textbf{Function} & \textbf{Reason} & \textbf{Requirements} \\
        \hline
        1 & SCROLL A & Close to where rotary encoder was placed & Requires an interrupt pin (all GPIOs on this processor can be set as interrupts) \\ \hline
        2 & OCTOSPIM P1 DQS & Auto assigned by STM32CubeIDE after selecting 'PORT1 DQS' Data Strobe & nil. \\ \hline
        3 & DCMI D4 & Auto assigned by STM32CubeIDE after selecting 'Slave 12 bits External Synchro' on DCMI \hypertarget{dcmi-reason}{[1]} & nil. \\ \hline
        4 & DCMI D6 & \hyperlink{dcmi-reason}{1} & nil. \\ \hline
        5 & DCMI D7 & \hyperlink{dcmi-reason}{1} & nil. \\ \hline
        6 & 3V3 & No external battery used & Tie to VCC if no external battery. \\ \hline
        7 & SCROLL B & Close to where rotary encoder was placed & Requires an interrupt pin (all GPIOs on this processor can be set as interrupts) \\ \hline
        8 & LSE IN & Input pin for low speed clock & 32.768kHz \\ \hline
        9 & NC & Output of LSE & Leave NC if using a single output oscillator (see Table 3 of getting started guide) \\ \hline
        10 & I2C2 SDA & Auto assigned by STM32CubeIDE after selecting 'I2C2' & nil. \\ \hline
        11 & I2C2 SCL & Auto assigned by STM32CubeIDE after selecting 'I2C2' & nil. \\ \hline
        12 & NC & nil. & nil. \\ \hline
        13 & 3V3 & Same as VCC & Must be connected to same supply as VCC \\ \hline
        14 & VDDCORE & Same voltage as VDDCORE & Must be connected to same supply as VDD11 \\ \hline
        15 & NC & nil. & nil. \\ \hline
        16 & NC & nil. & nil. \\ \hline
        17 & GND & VSS pin = GND & nil. \\ \hline
        18 & NC & nil. & nil. \\ \hline
        19 & NC & nil. & nil. \\ \hline
        20 & VDDCORE & Same voltage as VDDCORE & Must be connected to same supply as VDD11 \\ \hline
        21 & NC & nil. & nil. \\ \hline
        22 & NC & nil. & nil. \\ \hline
        23 & GND & VSS pin = GND & nil. \\ \hline
        24 & GND & VSS pin = GND & nil. \\ \hline
        25 & 3V3 & Common power domain & 1.71-3.6V \\ \hline
        26 & HSE IN & Input pin for HSE & 4-50MHz \\ \hline
        27 & NC & Output of HSE & Leave NC if using a single output oscillator (see Table 3 of getting started guide) \\ \hline
        28 & NRST & nil. & Active low \\ \hline
        29 & PWR ACTION & Close to MOSFET & nil. \\ \hline
        30 & PWR STATE & Close to MOSFET & nil. \\ \hline
    \end{tabular}
    \caption{Pin justification for STM32U5 (Pins 1-30).}
    \label{tab:stm32u5-pin-justification}
\end{table}

\begin{table}[H]
    \ContinuedFloat
    \centering
    \begin{tabular}{|l|p{3cm}|p{6.5cm}|p{3.5cm}|}
        \hline
        \textbf{Pin} & \textbf{Function} & \textbf{Reason} & \textbf{Requirements} \\
        \hline
        31 & CSA OUT & Close to Current Sense Amplifier & Must be analog GPIO \\ \hline
        32 & NC & nil. & nil. \\ \hline
        33 & GND & VSS pin = GND & nil. \\ \hline
        34 & 1V8 & Same supply as VDDA & Connect to VDDA \\ \hline
        35 & 3V3 & Common power domain & 1.62-3.6V \\ \hline
        36 & NC & nil. & nil. \\ \hline
        37 & SPI1 CLK & Auto assigned by STM32CubeIDE after selecting 'Half-Duplex Master' on Mode & nil. \\ \hline
        38 & USART2 TX & Auto assigned by STM32CubeIDE after selecting 'Asynchronous' on Mode & nil. \\ \hline
        39 & USART2 RX & Auto assigned by STM32CubeIDE after selecting 'Asynchronous' on Mode & nil. \\ \hline
        40 & GND & VSS pin = GND & nil. \\ \hline
        41 & 3V3 & Common power domain & 1.71-3.6V \\ \hline
        42 & DCMI HSYNC & \hyperlink{dcmi-reason}{1} & nil. \\ \hline
        43 & NC & nil. & nil. \\ \hline
        44 & DCMI CLK & \hyperlink{dcmi-reason}{1} & nil. \\ \hline
        45 & OCTIOSPIM P1 IO2 & Auto assigned by STM32CubeIDE after selecting 'PORT1 IO[3:0]' on Data [3:0] \hypertarget{octo-io-4+}{[2]} & . \\ \hline
        46 & OCTIOSPIM P1 IO1 & Auto assigned by STM32CubeIDE after selecting 'PORT1 IO[3:0]' on Data [3:0] \hypertarget{octo-io-4-}{[3]} & nil. \\ \hline
        47 & OCTIOSPIM P1 IO0 & \hyperlink{octo-io-4-}{3} & nil. \\ \hline
        48 & NC & nil. & nil. \\ \hline
        49 & NC & nil. & nil. \\ \hline
        50 & SENSOR RST TP & Near available space for test points & GPIO \\ \hline
        51 & NC & nil. & nil. \\ \hline
        52 & NC & nil. & nil. \\ \hline
        53 & TIM1 CH1 (Buzzer) & Auto assigned by STM32CubeIDE after selecting 'PWM Generation CH1' on Channel1 & nil. \\ \hline
        54 & GND & VSS pin = GND & nil. \\ \hline
        55 & 3V3 & Common power domain & 1.71-3.6V \\ \hline
        56 & OCTOSPIM P1 CLK & Auto assigned by STM32CubeIDE after selecting 'Port1 CLK' on Clock & nil. \\ \hline
        57 & OCTOSPIM P1 NCS & Auto assigned by STM32CubeIDE after selecting 'Port1 NCS' on 'Chip Select' & nil. \\ \hline
        58 & NC & nil. & nil. \\ \hline
        59 & SENSOR LED 2 & Near LED array & GPIO \\ \hline
        60 & SENSOR LED 1 & Near LED array & GPIO \\ \hline
        61 & OCTIOSPIM P1 IO3 & Auto assigned by STM32CubeIDE after selecting 'PORT1 IO[3:0]' on Data [3:0] \hypertarget{octo-io-4+}{[2]} & . \\ \hline
        62 & I2C4 SDA & Auto assigned by STM32CubeIDE after selecting 'I2C4' & nil. \\ \hline
        63 & I2C4 SCL & Auto assigned by STM32CubeIDE after selecting 'I2C4' & nil. \\ \hline
        64 & VDDCORE & Source of VDDCORE & nil. \\ \hline
        65 & 3V3 & Same source as VDD & Connect to VDD \\ \hline
        66 & GND & VSS pin = GND & nil. \\ \hline
        67 & VDDCORE & Same source as VLXSMPS & Connect to VLXSMPS \\ \hline
        68 & GND & VSS pin = GND & nil. \\ \hline
        69 & 3V3 & Common power domain & 1.71-3.6V \\ \hline
    \end{tabular}
    \caption{Pin justification for STM32U5 (continued, Pins 31-69).}
\end{table}

\begin{table}[H]
    \ContinuedFloat
    \centering
    \begin{tabular}{|l|p{3cm}|p{6.5cm}|p{3.5cm}|}
        \hline
        \textbf{Pin} & \textbf{Function} & \textbf{Reason} & \textbf{Requirements} \\
        \hline
        70 & NC & nil. & nil. \\ \hline
        71 & MCU LED 1 & Near LED array & GPIO \\ \hline
        71 & MCU LED 2 & Near LED array & GPIO \\ \hline
        73 & READ LED 1 & Near LED array & GPIO \\ \hline
        74 & READ LED 2 & Near LED array & GPIO \\ \hline
        75 & WRITE LED 1 & Near LED array & GPIO \\ \hline
        76 & WRITE LED 2 & Near LED array & GPIO \\ \hline
        77 & DEBUG LED & Near LED array & GPIO \\ \hline
        78 & NC & nil. & nil. \\ \hline
        79 & 3V3 LED & Near LED array & GPIO \\ \hline
        80 & NC & nil. & nil. \\ \hline
        81 & 81 & Originally a HSPI pin, repurposed as a broken out GPIO & GPIO \\ \hline
        82 & GND & VSS pin = GND & nil. \\ \hline
        83 & 3V3 & Common power domain & 1.71-3.6V \\ \hline
        84 & 84 & Originally a HSPI pin, repurposed as a broken out GPIO & GPIO \\ \hline
        85 & 85 & Originally a HSPI pin, repurposed as a broken out GPIO & GPIO \\ \hline
        86 & DCMI D3 & \hyperlink{dcmi-reason}{1} & nil. \\ \hline
        87 & NC & nil. & nil. \\ \hline
        88 & GND & VSS pin = GND & nil. \\ \hline
        89 & 3V3 & Common power domain & 1.71-3.6V \\ \hline
        90 & NC & nil. & nil. \\ \hline
        91 & DCMI D11 & \hyperlink{dcmi-reason}{1} & nil. \\ \hline
        92 & NC & nil. & nil. \\ \hline
        93 & DCMI D8 & \hyperlink{dcmi-reason}{1} & nil. \\ \hline
        94 & GND & VSS pin = GND & nil. \\ \hline
        95 & 3V3 & Common power domain & 1.71-3.6V \\ \hline
        96 & DCMI D9 & \hyperlink{dcmi-reason}{1} & nil. \\ \hline
        97 & 97 & Originally a HSPI pin, repurposed as a broken out GPIO & GPIO \\ \hline
        98 & DCMI D0 & \hyperlink{dcmi-reason}{1} & nil. \\ \hline
        99 & DCMI D1 & \hyperlink{dcmi-reason}{1} & nil. \\ \hline
        100 & SDMMC1 DAT0 & Auto assigned after selecting 'SD 4 bits Wide with auto dir voltage converter' on Mode \hypertarget{sdmmc1}{[5]} & nil. \\ \hline
        101 & SDMMC1 DAT1 & \hyperlink{sdmmc1}{5} & nil. \\ \hline
        102 & NC & nil. & nil. \\ \hline
        103 & USART1 TX & Auto assigned by STM32CubeIDE after selecting 'Asynchronous' on Mode & nil. \\ \hline
        104 & USART1 RX & Auto assigned by STM32CubeIDE after selecting 'Asynchronous' on Mode & nil. \\ \hline
        105 & NC & nil. & nil. \\ \hline
        106 & SPI1 MOSI & Auto assigned by STM32CubeIDE after selecting 'Half-Duplex Master' on Mode & nil. \\ \hline
        107 & SWDIO & Auto assigned by STM32CubeIDE after sleecting 'Serial Wire' on Debug & nil. \\ \hline
        108 & 3V3 & USB not used as a transceiver & Must be connected to VDD if transceiver not used \\ \hline
        109 & 3V3 & Common power domain & 1.71-3.6V \\ \hline
        110 & GND & VSS pin = GND & nil. \\ \hline
    \end{tabular}
    \caption{Pin justification for STM32U5 (continued, Pins 70-110).}
\end{table}

\begin{table}[H]
    \ContinuedFloat
    %\centering
    \begin{tabular}{|l|p{3cm}|p{6.5cm}|p{3.5cm}|}
        \hline
        \textbf{Pin} & \textbf{Function} & \textbf{Reason} & \textbf{Requirements} \\
        \hline
        111 & SWDCLK & Auto assigned by STM32CubeIDE after sleecting 'Serial Wire' on Debug & nil. \\ \hline
        112 & NC & nil. & nil. \\ \hline
        113 & SDMMC1 DAT2 & \hyperlink{sdmmc1}{5} & nil. \\ \hline
        114 & SDMMC1 DAT3 & \hyperlink{sdmmc1}{5} & nil. \\ \hline
        115 & SDMMC1 CLK & \hyperlink{sdmmc1}{5} & nil. \\ \hline
        116 & NC & nil. & nil. \\ \hline
        117 & NC & nil. & nil. \\ \hline
        118 & SDMMC1 CMD & \hyperlink{sdmmc1}{5} & nil. \\ \hline
        119 & DCMI D5 & \hyperlink{dcmi-reason}{1} & nil. \\ \hline
        120 & NC & nil. & nil. \\ \hline
        121 & NC & nil. & nil. \\ \hline
        122 & 3V3 & Common power domain & 1.71-3.6V \\ \hline
        123 & DMCI D10 & \hyperlink{dcmi-reason}{1} & nil. \\ \hline
        124 & NC & nil. & nil. \\ \hline
        125 & NC & nil. & nil. \\ \hline
        126 & NC & nil. & nil. \\ \hline
        127 & NC & nil. & nil. \\ \hline
        128 & I2C1 SDA & Auto assigned by STM32CubeIDE after selecting 'I2C1' & nil. \\ \hline
        129 & I2C1 SCL & Auto assigned by STM32CubeIDE after selecting 'I2C1' & nil. \\ \hline
        130 & GND & VSS pin = GND & nil. \\ \hline
        131 & 3V3 & PG[15:2] doesn't require a different power domain, I2C commonly runs on 3V3 logic & 1.08-3.6V. Independent from VDD when PG[15:2] in use, tie to VDD when not in use \\ \hline
        132 & STANDBY & Close to sensor connector & GPIO \\ \hline
        133 & RESET & Close to sensor connector & GPIO \\ \hline
        134 & TRIGGER & Close to sensor connector & GPIO \\ \hline
        135 & NC & nil. & nil. \\ \hline
        136 & SHUTTER & lose to sensor connector & GPIO \\ \hline
        137 & DCMI VSYNC & \hyperlink{dcmi-reason}{1} & nil. \\ \hline
        138 & BOOT0 & nil. & Active high \\ \hline
        139 & SDMMC1 CLKIN & \hyperlink{sdmmc1}{5} & nil. \\ \hline
        140 & NC & nil. & nil. \\ \hline
        141 & DCMI D2 & \hyperlink{dcmi-reason}{1} & nil. \\ \hline
        142 & VDDCORE & Same voltage as VDDCORE & Must be connected to same supply as VDD11 \\ \hline
        143 & SCROLL BUTTON & Close to where rotary encoder was placed & Requires an interrupt pin (all GPIOs on this processor can be set as interrupts) \\ \hline
        144 & 3V3 & Common power domain & 1.71-3.6V \\ \hline
    \end{tabular}
    \caption{Pin justification for STM32U5 (continued, Pins 111-144).}
    \label{tab:v1-pin-justification}
\end{table}

\subsection{PCB Layout}

\begin{indentedfigure}
    \includegraphics[width=0.6\linewidth]{Images/v0.1/PCB/length rule.png}
    \caption{FMC length match design rule.}
    \label{fig:length-rule}
\end{indentedfigure}

\begin{indentedfigure}
    \includegraphics[width=0.6\linewidth]{Images/v0.1/PCB/parallel routing.png}
    \caption{Parallel routes from C2Sv0.1 PCB.}
    \label{fig:parallel-routing}
\end{indentedfigure}
